On the one hand, by~1.2.4, \(f\) defines an element \(a\) of
\(C^{1}(Y_0,L_{0X})\) such that \(f(y)=a(y_0)y\), for every
\(y\in Y(S')\), \(S'\to S\), and by~4.27~(i), one has
\[
d^{1}(a)=d(Y,Y').
\]
Moreover, as \(Y,Y'\) are subgroups of \(X\), the above element belongs to
\(Z^{1}(Y_0,N_0)\) (\cf~4.21). Then \(\partial a\) is an element of
\(Z^{2}(Y_0,L_{0Y})\) whose image \(\partial a\) in \(H^{2}(Y_0,L_{0Y})\)
depends only on the class \(d(Y,Y')\in H^{1}(Y_0,N_0)\); this being the
definition of the boundary map
\(\partial^{1}\colon H^{1}(Y_0,N_0)\to H^{2}(Y_0,L_{0Y})\), one therefore
has:
\[
\partial^{1}\bigl(d(Y,Y')\bigr)=\partial a.
\]

On the other hand, let us transport by \(f\) the group structure of
\(Y'\) and let \(Y_1\) be the group obtained (which therefore has \(Y\) as
underlying scheme), that is, the group law \(\mu_1\) of \(Y_1\) is defined
by: for every \(S'\to S\) and \(x,y\in Y(S')\),
\[
\mu_1(x,y)=f^{-1}\bigl(f(x)f(y)\bigr).
\]
By~3.5.1, \(Y_1\) defines a cocycle
\(\delta(Y,Y_1)\in Z^{2}\bigl(Y_0,\mathcal{L}\mathrm{ie}(Y_0/S_0)\bigr)\)
such that, for every \(S'\to S\) and \(x,y\in Y(S')\), one has
\[
\delta(Y,Y_1)(x_0,y_0)\,xy=\mu_1(x,y)=f^{-1}\bigl(f(x)f(y)\bigr).
\]
Let us set \(b=\delta(Y,Y_1)\). For every \(S'\to S\) and
\(x,y\in Y(S')\), one has \((b(x_0,y_0)\,xy)_0=x_0y_0\), and hence one
obtains that \(f(b(x_0,y_0)\,xy)\) equals, on the one hand,
% typo?: kept as printed (\(b(x_0 y_0)\), without a comma)
\(a(x_0y_0)\,b(x_0y_0)\,xy\) and, on the other hand,
\[
f(x)f(y)=a(x_0)\,x\,a(y_0)\,y=a(x_0)\,\Ad(x_0)\bigl(a(y_0)\bigr)\,xy.
\]
Comparing the two expressions, one obtains that \(b(x_0,y_0)\) equals
\begin{align*}
a(x_0y_0)^{-1}a(x_0)\,\Ad(x_0)\bigl(a(y_0)\bigr)
&=\Ad(x_0)\bigl(a(y_0)\bigr)-a(x_0y_0)+a(x_0)\\
&=(\partial a)(x_0,y_0),
\end{align*}
\ie \(\delta(Y,Y_1)=\partial a\). One has therefore obtained the

\begin{proposition}{4.35.1}
\label{III.4.35.1}
\textsuperscript{(132)}
Under the hypotheses of~4.21, suppose in addition that \(S_0\) is affine
and that \(Y_{\mathcal{J}}\) is smooth over \(S_{\mathcal{J}}\) and
affine. Let \(Y,Y'\) be two subgroups of \(X\) lifting \(Y_{\mathcal{J}}\)
and flat (hence smooth) over \(S\), let \(f\colon Y\isomto Y'\) be an
isomorphism of \(S\)-schemes inducing the identity on \(Y_{\mathcal{J}}\),
and denote by \(Y_1\) the group obtained by transporting by \(f\) the group
structure of \(Y'\). Then one has
\[
\partial^{1}d(Y,Y')=\delta(Y,Y_1).
\]
\end{proposition}

\begin{proposition}{4.36}
\label{III.4.36}
Under the hypotheses of~4.21, suppose in addition that \(Y_{\mathcal{J}}\)
is smooth over \(S_{\mathcal{J}}\) and that \(S_0\) is affine. The set of
\(S\)-subgroups \(Y\) of \(X\) flat (or smooth) over \(S\), reducing to
\(Y_{\mathcal{J}}\), modulo conjugation by sections of \(X\) over \(S\)
inducing the unit section of \(X_{\mathcal{J}}\), is either empty or a
principal homogeneous set under the group
\[
H^{1}\bigl(Y_0,
\bigl[\mathcal{L}\mathrm{ie}(X_0/S_0)/\mathcal{L}\mathrm{ie}(Y_0/S_0)\bigr]
\otimes_{\OO_{S_0}}\mathcal{J}\bigr).
\]
\end{proposition}

It is enough for us to check that Corollary~4.29 applies, that is, that
\[
d^{0}\colon\Hom_{\OO_{S_0}}(\omega^{1}_{X_0/S_0},\mathcal{J})
\longrightarrow
\Hom_{\OO_{S_0}}(\mathfrak{n}_{Y_0/X_0},\mathcal{J})
\]
is surjective. Now this follows from the fact that the sequence~\((+)\)
of~4.25~(ii) is split, \(S_0\) being
affine.{\renewcommand{\thefootnote}{(133)}\nde{This also follows from the
proof of~4.32.}}

Let us finally state a corollary common to~4.21 and~4.36, which will, in
fact, be the only form in which we shall use in the sequel the general
results of this
number.{\renewcommand{\thefootnote}{(134)}\nde{For example, 4.37 is used
in Exposé~IX to prove the statements 3.2~bis and 3.6~bis.}}

\begin{corollary}{4.37}
\label{III.4.37}
Let \(S\) be a scheme and \(S_0\) the closed subscheme defined by a
nilpotent ideal \(\mathcal{I}\). Let \(X\) be an \(S\)-group smooth over
\(S\), and \(Y_0\) an \(S_0\)-subgroup of \(X_0\), flat over \(S_0\).
\begin{enumeratei}
\item If \(S_0\) is affine, \(Y_0\) is smooth over \(S_0\), and if
\[
H^{1}\Bigl(Y_0,
\bigl[\mathcal{L}\mathrm{ie}(X_0/S_0)/\mathcal{L}\mathrm{ie}(Y_0/S_0)\bigr]
\otimes_{\OO_{S_0}}\mathcal{I}^{n+1}/\mathcal{I}^{n+2}\Bigr)=0
\]
for every \(n\geqslant 0\), two \(S\)-subgroups of \(X\), flat (or smooth)
over \(S\), reducing to \(Y_0\), are conjugate by a section of \(X\) over
\(S\) inducing the unit section of \(X_0\).
\item If \(Y_0\) is affine and of finite presentation and
if{\renewcommand{\thefootnote}{(135)}\nde{One has replaced
\(\mathcal{H}\mathrm{om}_{\OO_{S_0}}(\mathfrak{n}_{Y_0/X_0},
\mathcal{I}^{n+1}/\mathcal{I}^{n+2})\) by
\(\mathfrak{n}^{\vee}_{Y_0/X_0}\otimes_{\OO_{S_0}}
\mathcal{I}^{n+1}/\mathcal{I}^{n+2}\), in accordance with Remark~4.22.}}
\[
H^{2}\bigl(Y_0,\mathfrak{n}^{\vee}_{Y_0/X_0}\otimes_{\OO_{S_0}}
\mathcal{I}^{n+1}/\mathcal{I}^{n+2}\bigr)=0
\]
for every \(n\geqslant 0\), there exists an \(S\)-subgroup of \(X\), flat
over \(S\), reducing to \(Y_0\).
\end{enumeratei}
\end{corollary}

\numpara{4.38}
\label{III.4.38}
It remains for us to relate the three constructions that we have made in
this exposé. In order to avoid inessential complications, we shall place
ourselves in the following situation: \(S_0\) is the spectrum of a field
\(k\), \(S\) is the spectrum of the dual numbers \(D(k)\), \(G\) is an
\(S\)-group smooth over \(S\), and \(K\) an \(S\)-subgroup, \emph{smooth}
over \(S\) and \emph{affine}.
{\renewcommand{\thefootnote}{(136)}\nde{One has slightly modified the
original in what follows. In particular, one has replaced \(X\) by \(G\)
and \(Y\) by \(K\), and one has denoted by \(\mathfrak{g}_0\) and
\(\mathfrak{k}_0\) their Lie algebras. On the other hand, one has written
explicitly \(H^{i}(K_0,\ )\) in place of the abbreviation \(H^{i}(\ )\) of
the original.}}
Let us denote \(\mathfrak{g}_0=\mathcal{L}\mathrm{ie}\,G_0\) (which here
equals \(\Gamma(S_0,\mathcal{L}\mathrm{ie}\,G_0)
=\Lie(G_0/S_0)(S_0)\)) and \(\mathfrak{k}_0=\mathcal{L}\mathrm{ie}\,K_0\).
One has an exact sequence of \(k\)-vector
spaces{\renewcommand{\thefootnote}{(137)}\nde{Equipped with the adjoint
action of \(K_0\).}}:
\[
0\longrightarrow\mathfrak{k}_0\xrightarrow{i}\mathfrak{g}_0
\xrightarrow{d}\mathfrak{n}^{\vee}_{K_0/G_0}\longrightarrow 0,
\]
giving rise to an exact cohomology sequence:
\begin{multline*}
0\longrightarrow H^{0}(K_0,\mathfrak{k}_0)
\xrightarrow{i^{0}}H^{0}(K_0,\mathfrak{g}_0)
\xrightarrow{d^{0}}H^{0}(K_0,\mathfrak{g}_0/\mathfrak{k}_0)\\
\xrightarrow{\partial^{0}}H^{1}(K_0,\mathfrak{k}_0)
\xrightarrow{i^{1}}H^{1}(K_0,\mathfrak{g}_0)
\xrightarrow{d^{1}}H^{1}(K_0,\mathfrak{n}^{\vee}_{K_0/G_0})
\xrightarrow{\partial^{1}}H^{2}(K_0,\mathfrak{k}_0).
\end{multline*}

Now these various groups all have a geometric meaning.

\noindent\textbf{a)}
\(H^{0}(K_0,\mathfrak{k}_0)
=\mathcal{L}\mathrm{ie}\,\Centr(K_0)\)%
{\renewcommand{\thefootnote}{(138)}\nde{As the formation of centralizers
and normalizers commutes with base change (\cf~I~2.3.3.1), one has written
\(\Centr(K_0)\) in place of \(\Centr(K)_0\) in the original, and likewise
\(\Centr_{G_0}(K_0)\) and \(\Norm_{G_0}(K_0)\) in place of
\(\Centr_{G}(K)_0\) and \(\Norm_{G}(K)_0\).}}
by~II~5.2.3.

\noindent\textbf{b)}
\(H^{0}(K_0,\mathfrak{g}_0)
=\mathcal{L}\mathrm{ie}\,\Centr_{G_0}(K_0)\)%
\textsuperscript{(138)}
(\emph{idem}).

\noindent\textbf{c)}
\(H^{0}(K_0,\mathfrak{g}_0/\mathfrak{k}_0)
=\mathcal{L}\mathrm{ie}\,\Norm_{G_0}(K_0)/\mathfrak{k}_0\)%
\textsuperscript{(138)}
(\emph{idem}).

\noindent\textbf{d)}
\(H^{1}(K_0,\mathfrak{k}_0)
=\mathcal{L}\mathrm{ie}\bigl(\Aut_{S_0\text{-gr.}}(K_0)\bigr)
/\Im(\mathfrak{k}_0)\),
where \(\Im(\mathfrak{k}_0)\) denotes the image of \(\mathfrak{k}_0\) under
the morphism \(\mathcal{L}\mathrm{ie}(\Int_0)\) deduced from
\(\Int_0\colon K_0\to\Aut_{S_0\text{-gr.}}(K_0)\). Indeed, it follows
from~2.1~(ii), applied to \(Y=X=K\) and \(f_0=\mathrm{id}_{K_0}\), that
\(Z^{1}(K_0,\mathfrak{k}_0)\) is the group of infinitesimal automorphisms of
the \(S_0\)-group \(K_0\), and that \(H^{1}(K_0,\mathfrak{k}_0)\) is
obtained by quotienting by the inner infinitesimal automorphisms, \ie by
the image of \(\mathfrak{k}_0\). Moreover, by~II~4.2.2, one also has
\(Z^{1}(K_0,\mathfrak{k}_0)
=\mathcal{L}\mathrm{ie}\bigl(\Aut_{S_0\text{-gr.}}(K_0)\bigr)\).%
{\renewcommand{\thefootnote}{(139)}\nde{And this is the Lie algebra
\(\operatorname{Der}_{k}(\mathfrak{k}_0)\) of the derivations of
\(\mathfrak{k}_0\), so \(H^{1}(K_0,\mathfrak{k}_0)\) is the quotient of
\(\operatorname{Der}_{k}(\mathfrak{k}_0)\) by the inner derivations
(\ie by the image of
\(\ad\colon\mathfrak{k}_0\to\operatorname{Der}_{k}(\mathfrak{k}_0)\)).}}

\noindent\textbf{e)}
\(H^{1}(K_0,\mathfrak{g}_0)\) is, by~2.1~(ii), the group of deviations
between homomorphisms \(K\to G\) extending the canonical immersion
\(i_0\colon K_0\to G_0\), modulo the deviations obtained by the action of
the inner automorphisms of \(G\) defined by elements of \(G(S)\) giving the
unit of \(G(S_0)\) (that is, elements of \(\mathfrak{g}_0\)).

\noindent\textbf{f)}
\(H^{1}(K_0,\mathfrak{n}^{\vee}_{K_0/G_0})\) is, by~4.36, the group of
deviations between subgroups \(K'\) of \(G\) extending \(K_0\) and flat over
\(S\) (hence \emph{smooth} over \(S\), \cf~SGA~1, II~4.10), modulo the
deviations obtained by the action of the inner automorphisms of \(G\)
constructed as previously.

\noindent\textbf{g)}
\(H^{2}(K_0,\mathfrak{k}_0)\) is, by~3.5~(ii), the group of deviations
between group structures on \(K\) extending that of \(K_0\), modulo the
\(S\)-automorphisms of \(K\) inducing the identity on \(K_0\).

We now propose to show how one can make explicit the six morphisms of the
preceding exact sequence in the geometric interpretation which we have just
given.

\noindent\textbf{1)}
\(i^{0}\) and \(d^{0}\) are none other than the morphisms obtained by
passage to the Lie algebra (then by passage to the quotient, for
\(d^{0}\)), starting from the canonical monomorphisms:
\[
\Centr(K_0)\longrightarrow\Centr_{G_0}(K_0)
\longrightarrow\Norm_{G_0}(K_0).
\]
This is indeed what follows immediately from the definition of the
identifications~(a), (b), and~(c).

\noindent\textbf{2)}
One constructs \(\partial^{0}\) as follows. Let
\(x\in\mathcal{L}\mathrm{ie}\,\Norm_{G_0}(K_0)/\mathfrak{k}_0\). Let us
lift it to an
\(x\in\mathcal{L}\mathrm{ie}\,\Norm_{G_0}(K_0)\subset\Norm_{G}(K)(S)\).
Then \(\Int(x)\) defines an automorphism of \(K\) inducing the identity on
\(K_0\), hence an element of
\(\mathcal{L}\mathrm{ie}\,\Aut_{S_0\text{-gr.}}(K_0)\). Denote by
\(\Int(x)\) the image of this element in
\(\mathcal{L}\mathrm{ie}\,\Aut_{S_0\text{-gr.}}(K_0)/\Im(\mathfrak{k}_0)\).
Then one has:
\[
\text{(*)}\qquad
\partial^{0}(x)=-\Int(x)=\Int(x^{-1}).
\]
Indeed, let us compute the element of
\(\mathcal{L}\mathrm{ie}\,\Aut_{S_0\text{-gr.}}(K_0)\) defined by
\(\Int(x)\). By definition it will correspond to an element \(a\) of
\(Z^{1}(K_0,\mathfrak{k}_0)\) such that
\[
xyx^{-1}=a(y_0)\,y,
\qquad\text{for every}\quad y\in K(S'),\ S'\longrightarrow S.
\]
But this may also be written
\(a(y_0)=xyx^{-1}y^{-1}=x-\Ad(y)x=-\partial(x)(y_0)\), whence
\(a=-\partial(x)\).

{\renewcommand{\thefootnote}{(140)}\nde{One has added the following
sentence.}}
On the other hand, the image of
\(x\in\mathcal{L}\mathrm{ie}\,\Norm_{G_0}(K_0)\subset\mathfrak{g}_0\) under
\(\partial\) is an element of \(Z^{1}(K_0,\mathfrak{k}_0)\), whose image
\(\partial(x)\) in \(H^{1}(K_0,\mathfrak{k}_0)\) depends only on \(x\), and
by definition of the boundary map \(\partial^{0}\), one has
\[
\partial^{0}(x)=\partial(x);
\]
combined with the equality \(a=-\partial(x)\), this proves~\((*)\).

\noindent\textbf{3)}
{\renewcommand{\thefootnote}{(141)}\nde{One has set out the original in
more detail in what follows, and in~\((**)\) one has corrected \(u\circ i\)
to \(i\circ u\).}}
Denote by \(i\colon K\to G\) the canonical immersion. Let \(u\) be an
element of \(H^{1}(K_0,\mathfrak{k}_0)\), the image of an
\[
u\in\mathcal{L}\mathrm{ie}\,\Aut_{S_0\text{-gr.}}(K_0)
\subset\Aut_{S\text{-gr.}}(K).
\]
Then one has:
\[
\text{(**)}\qquad i^{1}(u)=d(i,i\circ u),
\]
where \(d(i,i\circ u)\) is the class defined in~2.1.1.

Indeed, \(u\) is the image of an element \(v\in Z^{1}(K_0,\mathfrak{k}_0)\)
such that \(u(y)=v(y_0)y\), and \(i^{1}(u)\) is the image in
\(H^{1}(K_0,\mathfrak{g}_0)\) of the cocycle
% typo: cocyle -> cocycle
\(i\circ v\in Z^{1}(K_0,\mathfrak{g}_0)\).

Now, as \(i\) is a morphism of groups, the equality \(u(y)=v(y_0)y\)
implies \(iu(y)=iv(y_0)\,i(y)\). It follows that
\(i\circ v=d(i,i\circ u)\), whence~\((**)\).

\noindent\textbf{4)}\
Let \(i'\colon K\to G\) be a morphism of groups lifting \(i_0\), let
\(d(i,i')\) be the class defined in~2.1.1, and let
% typo?: kept as printed (\(\mathfrak{n}_{K_0/X_0}\), not \(\mathfrak{n}^{\vee}_{K_0/G_0}\))
\(d\bigl(K,i'(K)\bigr)\in C^{1}(K_0,\mathfrak{n}_{K_0/X_0})\) be the
deviation defined in~4.5.1; by~4.21, \(d\bigl(K,i'(K)\bigr)\) belongs to
\(Z^{1}(K_0,\mathfrak{n}_{K_0/X_0})\). Denote by \(d\bigl(K,i'(K)\bigr)\)
its image in \(H^{1}(K_0,\mathfrak{n}_{K_0/X_0})\). Then, by~4.27~(i), one
has:
\[
(\dagger)\qquad
d^{1}\bigl(d(i,i')\bigr)=d\bigl(K,i'(K)\bigr).
\]

\noindent\textbf{5)}\
Finally, let \(K'\) be a subgroup of \(G\) lifting \(K_0\) and flat, hence
\emph{smooth}, over \(S\). One has supposed that \(K_0\) is \emph{affine}.
Then one knows that \(K\) and \(K'\) are isomorphic as schemes extending
\(K_0\) (\cf~0.15), hence that there exists an isomorphism of \(S\)-schemes
\[
f\colon K\isomto K'
\]
inducing the identity on \(K_0\). Let us transport by \(f\) the group
structure of \(K'\) and let \(K_1\) be the group obtained (which therefore
has \(K\) as underlying scheme), that is, the group law \(\mu_1\) of
\(K_1\) is defined by: for every \(S'\to S\) and \(x,y\in K(S')\),
\[
\mu_1(x,y)=f^{-1}\bigl(f(x)f(y)\bigr).
\]
{\renewcommand{\thefootnote}{(142)}\nde{One has modified the original in
what follows, taking into account the additions made in~3.5.1 and~4.35.1.}}
By~3.5.1, \(K_1\) defines a cocycle
\(\delta(K,K_1)\in Z^{2}(K_0,\mathfrak{k}_0)\) such that, for every
\(S'\to S\) and \(x,y\in K(S')\), one has
\[
\delta(K,K_1)(x_0,y_0)\,xy=\mu_1(x,y)=f^{-1}\bigl(f(x)f(y)\bigr).
\]
Then, by~4.35.1, one has:
\[
(\ddagger)\qquad
\partial^{1}d(K,K')=\delta(K,K_1).
\]

\makeatletter
\renewcommand{\@bibtitlestyle}{%
  \par\addvspace{1.1\baselineskip}%
  \noindent{\large\bfseries Bibliography}%
  {\renewcommand{\thefootnote}{(143)}\nde{: additional references cited in this Exposé.}}%
  \par\addvspace{0.45\baselineskip}%
}
\makeatother
\begin{thebibliography}{DG70}

\bibitem[BAlg]{BAlg}
N.~Bourbaki, \emph{Algèbre}, Chap.~I--III, Hermann, 1970.

\bibitem[BAC]{BAC}
N.~Bourbaki, \emph{Algèbre commutative}, Chap.~I--IV, Masson, 1985.

\bibitem[DG70]{DG70}
M.~Demazure, P.~Gabriel, \emph{Groupes algébriques}, Masson \&
North-Holland, 1970.

\bibitem[Fr64]{Fr64}
P.~Freyd, \emph{Abelian categories}, Harper and Row, 1964.

\end{thebibliography}
